% -----------------------------------------------------------------------------------------------------
% Hypotheses A3 and A2. Statements restated from theory/A3_lkatyusha_lyapunov.tex (\restate); numbers quoted:
%   A3_summary.json (3 seeds; hard instance, logistic kappa = 1e3, 1e4): speedup_exact_over_lemma 4.668, 2.775, 4.117;
%     speedup_exact_over_rho 9.33, 3.70, 8.23; contraction_componentwise_max < 0; q_plus_frac_ge_1 0.0986, 0.0519,
%     0.1929.
%   A2_summary.json (one seed, kappa = 1e4): theta2 = 0 diverged on (H) (m = 64); stable anchor-free schemes 8.1-8.5
%     (logistic) and 20.1 ((H)) times fewer gradients than L-SVRG (A3 note, Outcome: "8 to 20 times").
%   A4_summary.json (3 seeds, kappa >= 10n): slopes of the re-tuned naive scheme 0.806 ((H)) and 0.814 (random spectra).
% -----------------------------------------------------------------------------------------------------
\subsection{L-Katyusha: Lyapunov proof, negative momentum, and naive acceleration}\label{sec:th-A3}

Hypothesis A3 asks where the negative momentum of loopless Katyusha (L-Katyusha) \cite{kovalev2020dont} enters a
Lyapunov proof written in the style of the book. For hypothesis A2, whether naive acceleration fails (Nesterov's
momentum applied to the loopless SVRG estimator, without negative momentum), we ask why the same proof does not close in
this case. Appendix~\ref{sec:lkatyusha} rewrites the proof of \cite{kovalev2020dont} in the notation of the book, under
the following assumption.
\restate{ass:lk}
In this subsection $L=R^2$ is the per-sample smoothness constant, not the smoothness constant $L_F$ of $F$, which
Section~\ref{sec:th-B2} denotes $L$.
L-Katyusha keeps a gradient sequence $y_t$, a mirror sequence $z_t$ and a reference point $w_t$ at which the full
gradient $\Fp(w_t)$ is known. We keep the notation of \cite{kovalev2020dont} for its parameters: $\theta_1,\theta_2$ are
momentum weights, not iterates, $\eta$ is a step parameter, not the minimizer $\etastar$, and $z_t$ is the mirror
sequence, not the SAGA memory. Let $\theta_1,\theta_2>0$ with $\theta_1+\theta_2\le1$, $p\in(0,1]$, and
\begin{equation*}
\eta:=\frac{\theta_2}{(1+\theta_2)\theta_1}.\tag{\ref*{eq:lk-eta}}
\end{equation*}
From $y_0=z_0=w_0=\theta_0$, with $\Fp(w_0)$ computed, an iteration reads
\begin{align*}
x_t &= \theta_1 z_t+\theta_2 w_t+(1-\theta_1-\theta_2)y_t, \tag{K1}\\
g_t &= f'_{i(t)}(x_t)-f'_{i(t)}(w_t)+\Fp(w_t), \tag{K2}\\
z_{t+1} &= \frac1{1+\eta\sigma}\Big(\eta\sigma\,x_t+z_t-\frac\eta L\,g_t\Big), \tag{K3}\\
y_{t+1} &= x_t+\theta_1(z_{t+1}-z_t), \tag{K4}\\
w_{t+1} &= y_t\ \text{with probability $p$},\qquad w_{t+1}=w_t\ \text{otherwise}, \tag{K5}
\end{align*}
where $i(t)$ is uniform on $\{1,\dots,n\}$ and the coin of \eqref{eq:K5} is independent of $i(t)$ and of the past;
$\Ei$ denotes the expectation over both given $(y_t,z_t,w_t)$, so that $\Ei[g_t]=\Fp(x_t)$. Step \eqref{eq:K3} is a
mirror step, \eqref{eq:K4} a gradient step, and \eqref{eq:K1} couples them \cite{allenzhu2017linear}. The term
$\theta_2w_t$ of \eqref{eq:K1}, the anchor, is the \emph{negative momentum} of Katyusha \cite{allenzhu2018katyusha}: it
pulls the query point towards the point where the estimator has no variance. With $\theta_2=0$ and a free step $\eta$,
$x_{t+1}=y_{t+1}+(1-\theta_1)(y_{t+1}-y_t)$ and the method is a variant of Nesterov's momentum applied to the loopless
SVRG estimator, the naive acceleration of hypothesis A2; see \resref{rem:naive-book}. Theorem~11
of \cite{kovalev2020dont} takes
\begin{equation*}
\theta_2=\tfrac12,\qquad \theta_1=\min\Big\{\sqrt{\tfrac{2\sigma n}{3}},\,\tfrac12\Big\}=\min\Big\{\sqrt{\tfrac{2n}{3\kappa}},\,\tfrac12\Big\},\qquad p=\tfrac1n,\qquad \eta=\frac{\theta_2}{(1+\theta_2)\theta_1}=\frac1{3\theta_1},\tag{\ref*{eq:lk-param}}
\end{equation*}
and its proof tracks the Lyapunov function
\begin{equation*}
\begin{aligned}
\Psi_t &:= Z_t+Y_t+W_t,\\
Z_t&:=\frac{L(1+\eta\sigma)}{2\eta}\norm{z_t-\etastar}^2,\qquad Y_t:=\frac1{\theta_1}\big(F(y_t)-F_\ast\big),\qquad W_t:=\frac{\theta_2(1+\theta_1)}{p\theta_1}\big(F(w_t)-F_\ast\big),
\end{aligned}\tag{\ref*{eq:lk-lyap}}
\end{equation*}
the sum of a distance of the mirror sequence, of a function value of the gradient sequence and of a memory term, which
measures the quality of the reference point. Finally, $D_F(w,x):=F(w)-F(x)-\inner{\Fp(x)}{w-x}\ge0$ denotes the Bregman
divergence of $F$.

The proof certifies a contraction of $\Psi$ in conditional expectation at every step, for all admissible parameters,
and a global rate for the parameters \eqref{eq:lk-param}.
\restate{thm:lk}
We write $c_Z$, $c_Y$ and $c_W$ for the three coefficients of \eqref{ineq:contraction}, so that $\Ei[\Psi_{t+1}]$ is
at most $\max\{c_Z,c_Y,c_W\}\,\Psi_t$, a factor that the proof of \eqref{ineq:rate} bounds by $1-\rho$. The two
terms of $\rho$ are lower bounds on $1-c_Z$, the contraction rate of the distance $Z_t$ of the mirror sequence, whose
step $\eta/L=1/(3\theta_1L)$ grows as $\theta_1$ decreases, and on $1-c_W$, that of the memory term $W_t$, refreshed with
probability $\frac1n$; for $\kappa\ge\frac{8n}3$, the value $\theta_1=\sqrt{2n/(3\kappa)}$ of \eqref{eq:lk-param} balances
them up to a factor $2$. Along L-Katyusha runs, with
$\Ei[\Psi_{t+1}]$ computed exactly, \eqref{ineq:contraction} holds at every recorded state, and the exact one-step
contraction exponent is on average $2.8$ to $4.7$ times the certified exponent $-\log\max\{c_Z,c_Y,c_W\}$ and $3.7$ to
$9.3$ times $-\log(1-\rho)$ (three problems, 3 seeds; Section~\ref{sec:exp-A3}). In gradient evaluations, the rate gives
the following count.
\restate{cor:lk-complexity}
The count without storage, three gradients per iteration in expectation, is the convention of the experiments and of
Table~\ref{tab:th-complexity}; the count with storage is that of the bounds of Appendix~\ref{sec:lower_bound}. The
condition on $\varepsilon$ lets the constants absorb the rounding of the number of iterations to an integer. For
$\kappa\gg n$, the leading term $\sqrt{n\kappa}\log\frac1\varepsilon$ is of the order of the lower bounds of
Section~\ref{sec:th-A1}.

\paragraph{The proof, and where the negative momentum acts.}
The proof of the one-step contraction \eqref{ineq:contraction} chains five lemmas, L1--L5, which reproduce Lemmas~6--10
of \cite{kovalev2020dont} and use seven elementary inequalities, listed with their tightness ratios in
Table~\ref{tab:lk-instrument} of Appendix~\ref{sec:lkatyusha}. L1 (\resref*{lem:L1}) bounds the variance of the
estimator; L2 (\resref*{lem:L2}) bounds the mirror step by the three-point identity; L3 (\resref*{lem:L3}) bounds the
gradient step by smoothness and Young's inequality; L4 (\resref*{lem:L4}) is the identity \eqref{eq:wstep} for the
refresh of $w$; and L5 (\resref*{lem:L5}) combines them with strong convexity at $x_t$ and convexity towards $y_t$. The
form of L1 matters for hypothesis A3: it bounds the variance by a Bregman divergence, whose inner-product part the
negative momentum cancels.
\restate{lem:L1}
It is the one-sided (Bregman) form of co-coercivity in \eqref{ineq:var-coco}, instead of the symmetric form used in the
proof of Proposition~\ref{prop:59}, that produces this divergence; Section~\ref{sec:exp-A3} measures the two steps of
\eqref{ineq:var} along trajectories. In the proof of L5 (\resref*{lem:L5}), $\theta_2$ is used for something other
than bookkeeping at a single step, Step~6, the identity
\begin{equation*}
\frac{\theta_2}{\theta_1}\inner{\Fp(x_t)}{x_t-w_t}-\frac{\theta_2}{\theta_1}\Big(F(w_t)-F(x_t)-\inner{\Fp(x_t)}{w_t-x_t}\Big)=-\frac{\theta_2}{\theta_1}\big(F(w_t)-F(x_t)\big),\tag{\ref*{eq:cancel}}
\end{equation*}
in which the first term comes from the anchor $\theta_2w_t$ of \eqref{eq:K1} and the second is the variance term, which
enters through L3 and L1:
\restate{prop:cancel}
The inner-product part of $D_F(w_t,x_t)$ must be cancelled rather than absorbed. By smoothness,
$D_F(w_t,x_t)\le\frac L2\norm{w_t-x_t}^2$ is a squared distance at the scale $L$, while the certified decrease per step
is of order $\rho\Psi_t$, with $\rho=1/\sqrt{24n\kappa}$ when $\kappa\ge\frac{8n}3$: absorbing this term would require
converting a squared distance into function values, at the price of a factor $\kappa$, which \resref{prop:naive-kappa}
makes precise for the naive scheme; and the inner product $\inner{\Fp(x_t)}{x_t-w_t}$ has no sign and is not a
difference of function values that a potential could pay. The negative momentum costs a constant factor: $Y_t$
contracts by $\theta_1(1-\theta_2)$ instead of $\theta_1$, half as fast for $\theta_2=\frac12$. It also ties the step to
$\theta_2$, as \resref{rem:eta-theta2} explains: for a general step, Young's inequality in L3 (\resref*{lem:L3})
produces the variance coefficient $\nu/(2L\theta_1)$ with $\nu=\eta\theta_1/(1-\eta\theta_1)$, and the cancellation is
exact only for $\nu=\theta_2$, that is, for the step \eqref{eq:lk-eta}: a larger step leaves the residual
$\frac{\nu-\theta_2}{\theta_1}D_F(w_t,x_t)$, the extra term of the naive scheme (\resref*{prop:naive}) with
$\nu-\theta_2$ in place of $\nu$, and a smaller one wastes part of the pull. Along the L-Katyusha runs of
Section~\ref{sec:exp-A3}, the cancelled term $\frac{\theta_2}{\theta_1}\inner{\Fp(x_t)}{x_t-w_t}$ is positive and larger
than the certified one-step decrease $(1-\max\{c_Z,c_Y,c_W\})\Psi_t$ on $9.9\%$, $5.2\%$ and $19.3\%$ of the recorded
states of the three problems; only there could it not have been absorbed into that decrease, since where it is
negative, $D_F(w_t,x_t)\le F(w_t)-F(x_t)$ is bounded by a difference of function values.

\paragraph{Naive acceleration (hypothesis A2).}
Without the anchor, the naive scheme is the same iteration with $\theta_2=0$ and a free step.
\restate{def:naive}
With exact gradients, the same lemmas prove acceleration, so that what fails with the estimator is the treatment of
its variance.
\restate{prop:agd-lc}
With the estimator, Steps~1--5 and~7 of the proof of L5 (\resref*{lem:L5}) go through, but no step produces the inner
product that cancelled the variance term.
\restate{prop:naive}
The decrease of the potential that \eqref{ineq:naive} can certify per step is at most
$\frac{\eta\sigma}{1+\eta\sigma}Z_t+\theta_1Y_t$, plus, with a reference term $a_W(F(w_t)-F_\ast)$, the expected
decrease $pa_W(F(w_t)-F_\ast)-pa_W(F(y_t)-F_\ast)$ of that term; none of these involves the query point $x_t$.
Bounding the extra term through smoothness and strong convexity costs a factor $\kappa$:
\restate{prop:naive-kappa}
This is a limit of the Lyapunov argument, not a lower bound for the naive scheme, as \resref{rem:naive-honest}
details: bounding the extra term through $Z_t$ instead of $\mu$ is no better, paying it by the reference term alone is
impossible, and no lower bound specific to the scheme is known to us. It agrees with the literature:
with one sample per iteration, the guarantee of Nesterov's momentum on the SVRG estimator is that of SVRG
\cite{nitanda2014stochastic}, the anchor-free acceleration of \cite{driggs2022accelerating} needs mini-batches of size
of order $n^{2/3}$, and the semidefinite certificate of \cite{hu2018dissipativity} is stated for negative-momentum
weights of at least $\frac15$. In the experiments at $\kappa=10^4$ (one seed; Section~\ref{sec:exp-A2}), L-Katyusha
with $\theta_2=0$ and otherwise unchanged parameters diverges on the instance (H) ($n=16$ blocks of length $64$) and
converges on logistic regression, and Nesterov's momentum applied to the estimator, at its best grid setting, converges
on both problems; where they converge, these anchor-free schemes need $8$ to $20$ times fewer gradient evaluations
than L-SVRG. Re-tuned at every $\kappa$, Nesterov's momentum applied to the estimator has fitted slopes of about $0.8$
in $\kappa$ on block quadratics (Section~\ref{sec:exp-A4}). Prediction (P3) of Section~\ref{sec:lk-naive}, that a
stable naive scheme keeps the exponent $1$ of $\kappa$ and does no better than L-SVRG, is thus refuted in its second
part and not supported in its first: \resref*{prop:naive} and \resref*{prop:naive-kappa} describe a limitation of the
certificate, not the observed behaviour.

\paragraph{Known and new.}
\resref*{thm:lk} and its constants are those of \cite{kovalev2020dont}, the negative momentum is that of
\cite{allenzhu2018katyusha}, and its role of cancelling an inner product of the variance bound is stated in
\cite{allenzhu2018katyusha,zhou2019direct}. To our knowledge, the following are not in the sources we reviewed in this
form: the rewriting of the proof as a chain of labelled inequalities with tightness ratios, the isolation of the
cancellation as a single identity tied to the choice of the step (\resref*{prop:cancel} and \resref*{rem:eta-theta2}),
and \resref*{prop:naive} and \resref*{prop:naive-kappa}, which quantify what the same argument certifies without the
anchor.
